\chapter{预备知识}
本章基于在密文上的可验证查询框架，梳理出了在方案设计和具体实现中所用到的关键技术。
\section{可搜索加密技术}
数据外包可能会造成云服务商未经授权访问用户隐私，因此需要将本地文件加密后上传，但这会破坏数据的多种语义关系，从而导致传统搜索技术失效。因此，Song\}等人\cite{song2000practical}提出了可搜索加密技术。
它允许数据所有者将加密数据与索引结构存至云端，使得服务器可以仅凭借搜索令牌执行查询，返回密文文件。利用可搜索加密技术进行搜索，可以在节省用户本地存储空间的同时，保证了数据的机密性和隐私性。
这里接着给出可搜索加密算法和框架图。

可搜索加密技术基于加密算法可分为SSE技术和PEKS技术两大类。SSE包含四个信任实体：云服务商CSP、数据拥有者DO以及数据用户DU。DO一般会在本地生成密钥，负责将自己所拥有的共享数据加密上传给云服务商，并负责为DU进行密钥传送。
CSP负责为DO提供数据存储服务、为DU提供检索数据服务。DU现在本地利用自己所掌握的密钥信息，生成搜索陷门，并将相对应的搜索陷门提交给CSP，CSP会进行检索操作，返回DU所需的数据。PEKS技术也包含三个实体：DO、DU和CSP，
其中密钥由公钥和私钥组成，主要由可信结构生成。数据使用者DO负责将自己是所拥有的共享数据利用数据用户DU的公钥加密，然后上传给CSP。DU在本地使用私钥生成搜索陷门交给CSP，拿到返回的密文数据后，再利用自己私钥对数据进行解密。

然而无论是SSE还是PEKS技术，都存在xxxxx问题。在此基础上，发展出了基于区块链的可搜索加密技术，实体新增了一个“区块链。DO加密文件上传云服务器，另一方面构建加密索引上传之区块链。一般来说DO会创建智能合约，
并将加密索引存储在只能合约中，将智能合约部署在区块链平台。DU向数据所有者请求搜索某个关键字的权限。Owner验证User ，如果通过，则生成关键字的陷门并发送给DU。User 将搜索令牌发送给区块链上的智能合约。智能合约使用Token在加密索引中查找，返回对应的加密文件位置。
DU使用私钥解密得到明文文件位置，向云服务器请求该文件位置的加密文件。云服务器返回加密文件。DU使用对应密钥解密加密文件。
这个基于区块链的流程有些不太对，要给出区块链的系统模型图还有算法。

\section{哈希函数与伪随机函数}